\section{Background}\label{sec:background}

Throughout the paper we follow the notation in \autoref{fig:notation}.
%This section introduces the concepts in the figure.
We define the
 \egraph\ data structure and the \ematching\ problem,
 and review background on
 relational queries and join algorithms
 that form the foundation of our \ematching\ algorithm.

\subsection{E-Graphs and E-Matching}\label{sec:background:ematching}

\newcommand\nt[1]{\ensuremath{\langle\mathit{#1}\rangle}}
\begin{figure}
  \centering
  \begin{tabular}{llcl}
     function symbols & \nt{fun} & ::= & $f \mid g \mid \ldots$                                           \\
     variables        & \nt{var} & ::= & $x \mid y \mid z \mid \ldots \mid \alpha \mid \beta \mid \ldots$ \\
     \eclass ids      & \nt{id}  & ::= & $i \mid j \mid \ldots$                                           \\
     ground terms     & \nt{t}   & ::= & $\nt{fun} \mid \nt{fun}\left(\nt{t},\ldots,\nt{t}\right)$        \\
     patterns         & \nt{p}   & ::= & $\nt{fun} \mid \nt{fun}(\nt{p},\ldots,\nt{p}) \mid \nt{var}$     \\
     \enodes          & \nt{n}   & ::= & $\nt{fun} \mid \nt{fun}(\nt{id}, \ldots, \nt{id})$               \\
     \eclasses        & \nt{c}   & ::= & $\{ \nt{n}, \ldots, \nt{n}\}$                                    \\
  \end{tabular}
  \caption{Syntax and metavariables used in this paper.}
  \label{fig:notation}
\end{figure}

Let $\Sigma$ be a set of function symbols with associated arities. A function
symbol is called a \textit{constant} if it has zero arity. Let $V$ be a set
of variables. We define $T(\Sigma, V)$ to be the set of terms constructed using
function symbols from $\Sigma$ and variables from $V$:

\update{
\begin{definition}[Terms and patterns]
  The set of terms $T(\Sigma,V)$ over function symbols $\Sigma$ and variables $V$
 is the smallest set such that (1) all variables and constants are in
$T(\Sigma,V)$ and (2) $t_1,\dots,t_k\in T(\Sigma,V)$ implies
$f(t_1,\dots,t_k)\in T(\Sigma,V)$, where $f\in \Sigma$ has arity $k$. A
\textit{ground term} is a term in $T(\Sigma,V)$ that contains no variables.
All terms in $T(\Sigma, \emptyset)$ are ground terms.
A non-ground term is also called a \textit{pattern}. We call a term of the form
$f(t_1,\ldots,t_k)$ an $f$-application term.
\end{definition}


We define a {\em congruence relation} over the terms as an equivalence relation 
 that is {\em congruent}:

\begin{definition}[Equivalence Relation]
  An \textit{equivalence relation} $\equiv_\Sigma$ is a binary relation over
$T(\Sigma,\emptyset)$ that is reflexive, symmetric, and transitive. 
\end{definition}

\begin{definition}[Congruence Relation]
  A \textit{congruence relation} $\cong_\Sigma$ is an equivalence relation
satisfying:
\[
  \forall \text{$k$-ary function symbols $f$}.\
  \left(
    \forall i\in\{1, \ldots, k\}.\ t_i \cong t'_i
  \right)
  \implies
  f(t_1,\dots,t_k)\cong f(t'_1,\dots,t'_k)
\]  
The {\em congruence closure} of a binary relation $R$ on $T(\Sigma, \emptyset)$
 is the smallest congruence relation that contains $R$.
\end{definition}


We write $\equiv$ and $\cong$ when $\Sigma$ is clear from the context.

\newcommand{\child}{\textit{child}} \newcommand{\sym}{\textit{symbol}}

An \egraph is a data structure that represents a congruence relation.
An \egraph is built up from {\em \eclasses} and {\em \enodes}, defined as follows:
\begin{definition}[\Eclasses and \Enodes]
  An \eclass is a set of \enodes. Every \eclass is identified by one or more ids.
  An \enode is a tuple $(f, \text{args})$ where $f$ is a function symbol 
   and args is a (possibly empty) list of \eclass ids. 
Similar to terms, we call an \enode\ of the form $(f, i_1, \ldots i_k)$
 an $f$-application \enode.
We will write $f(i_1, \ldots, i_k)$ for the \enode $(f, i_1, \ldots i_k)$
\end{definition}

\begin{definition}
  A \egraph is a tuple $(C, E, I, U, M, \lookup)$ where:
  \begin{itemize}
    \item $C$ is a set of \eclasses over $E$ which is a set of \enodes;
    $I$ is a set of \eclass ids.
    \item
      A union-find~\cite{unionfind} data structure $U$
      stores an equivalence relation (denoted with $\equivid$)
      over \eclass ids.
      The union-find provides a function \find that
       \textit{canonicalizes} \eclass ids such that
       $\find(i_1) = \find(i_2) \iff i_1 \equivid i_2$.
      An \eclass id $i$ is canonical if $i = \find(i)$.
    
    \item
      The \textit{\eclass map} $M$ is a surjective function 
      that maps \eclass ids to \eclasses.
      All equivalent \eclass ids map to the same \eclass, i.e.,
      $a \equivid b$ iff $M[a]$ is the same set as $M[b]$.
      An \eclass id $a$ is said to \textit{refer to} the \eclass $M[\find(a)]$.
    
    \item
      A function \lookup that maps \enode $n$ to an id of
        \eclass that contains it: $n \in M[\lookup(n)]$.
    \end{itemize}
\end{definition}
}
 
Note that by definition, no two \enodes in the same \egraph 
 can have the same symbol and children, i.e.,
 an \enode's symbol and children together uniquely identify
 the \eclass that contains it.
This property is necessary for \lookup to be a function.
\autoref{sec:fd} explains how this property also
 translates to a {\em functional dependency}
 in the \egraph's relational representation,
 which could be leveraged to further optimize relational \ematching.

 \update{
\autoref{fig:egraph-eg} shows an example \egraph, where each dotted box
 is an \eclass, and each solid box together with argument pointers 
 make up an \enode. 
In this \egraph there is one unique id for each \eclass,
 shown in the shaded labels on the top-left corner of each \eclass. 
For example, the \eclass in the middle contains the set of \enodes 
$g(1), \ldots, g(N)$ and has id $i_g$. }

An \egraph $E$ efficiently represents sets of ground terms in a congruence
 relation.
An \egraph, \eclass, or \enode is said to \textit{represent} a term $t$ if $t$ can be
``found'' within it. 

\begin{definition}[Representation]
  Representation is defined recursively:
\begin{itemize}
\item An \egraph\ represents a term if any of its \eclasses\ does.
\item An \eclass\ $c$ represents a term if any \enode $n \in c$ does.
\item An \enode $f(j_{1}, \ldots, j_k)$ represents a term $f(t_{1}, \ldots, t_k)$
  if they have the same function symbol $f$
  and \eclass $M[j_{i}]$ represents term $t_{i}$ for $i \in \{1, \ldots, k\}$.
\end{itemize}
\end{definition}

%  An \egraph~is said to \textit{represent} a ground term $t$ if its
% \eclasses~represent it. An \eclass $c$ represents a ground term if $c$ contains
% an \enode $a$ that represents it. An \enode $f(c_1,\dots,c_k)$ represents a
% ground term $f(t_1,\dots,t_k)$ if they have the same function symbol $f$ and
% each e-class $c_i$ represents term $t_i$. Terms represented by an \eclass~are
% equivalent to each other.

\begin{figure}
  \centering
  \includegraphics[height=5cm]{fgn.pdf}
  \caption{An example \egraph.}\label{fig:egraph-eg}
\end{figure}

The \egraph in~\autoref{fig:egraph-eg} represents the set of
terms (where $[N] = \{1, 2, \ldots, N\}$):
$$[N]\cup \{g(i) \mid i\in [N]\} \cup\{f(i, g(j)) \mid i, j \in [N]\}.$$
In addition, all $g$-terms are equivalent,
 and all $f$-terms are equivalent.
Note that the \egraph~has size $O(N)$,
 yet it represents $\Omega(N^{2})$ many terms.
In general, an \egraph is capable of representing
 exponentially many terms in polynomial space.
If the \egraph has cycles,
 it can even represent an infinite set of terms.
For example, the \egraph with a single \eclass $c = \{f(c), a\}$
 represents the infinite set of terms $\{a, f(a), f(f(a)), \ldots\}$.


\subsubsection*{\Ematching.}
\Ematching~finds the set of terms in an \egraph~matching a given pattern.
Specifically, \ematching~finds the set of \ematching~substitutions and a root
class that represents the terms.

\begin{definition}[\Ematching substitution]
  An \textit{\ematching~substitution} $\sigma$ is a function that maps
every variable in a pattern to an \eclass.  
\end{definition}
 For convenience, we use $\sigma(p)$ to
denote the set of terms obtained by replacing every occurrence of variable $v_i$
in $p$ with terms represented by $\sigma(v_i)$.

\begin{definition}[The \Ematching problem]
  Given an e-graph $E$ and a pattern $p$,
\ematching~finds the set of all possible pairs $(\sigma, r)$ such that every
term in $\sigma(p)$ is represented in the \eclass $r$. Terms in $\sigma(p)$ are
said to be matched by pattern~$p$, and $r$ is said to be the root of matched terms.  
\end{definition}
For example, matching the pattern $f(\alpha, g(\alpha))$ against the \egraph~$G$
in \autoref{fig:egraph-eg} produces the following $N$ substitutions,
 each with the same root $i_f$:
\[
  \left\{
    (\{\alpha \mapsto j\}, i_f) \mid j \in [N]
  \right\}.
\]

Existing \ematching~algorithms perform backtracking search directly on the
\egraph~\cite{efficient-ematching,simplify,egg}.
\autoref{alg:backtrack} shows
 an abstract backtracking-based
 \ematching~algorithm.
Most
\ematching~algorithms using backtracking search can be viewed as
optimizations based on this abstract algorithm. Specifically,
it will perform a top-down search following the shape of the pattern and prune
the result set of substitutions when necessary.
To match pattern $f(\alpha,
g(\alpha))$ against $G$, backtracking search visits terms in the following order
(each $\hookrightarrow$ marks a backtrack step):
% TODO: Say something like f(1, g(1)) is not actually a term, but a (singleton)
% set of terms witnessed by the substitution {alpha -> 1}.
\begin{align*}
  f(1, g(1)) \rightarrow & \dots \rightarrow f(1, g(N)) \\
    \hookrightarrow f(2, g(1))\rightarrow & \dots \rightarrow f(2, g(N)) \\
    \hookrightarrow f(N, g(1))\rightarrow & \dots \rightarrow f(N, g(N))
\end{align*}

For each term $f(i, g(j))$ visited, whenever $i=j$ the algorithm yields a match
$\alpha \mapsto i$. Despite there being only $N$ matches,
backtracking search runs in time $O(N^{2})$.
% Existing approaches to \ematching~rely on backtracking. For example,
% \citet{efficient-ematching} proposed a backtracking-based e-matching algorithm
% that is used by \textsc{Z3} \cite{Z3} and \egg \cite{egg}, two
% state-of-the-art e-graph implementations. To match the pattern $f(\alpha,
% g(\alpha))$ on the e-graph in Figure \ref{fig:egraph}, their algorithm does a
% depth-first search over the e-graph:
% % Conceptually, it recursively searches for the pattern $f(\alpha, g(\beta))$
% % and filters out substitutions that map $\alpha$ and $\beta$ to different
% % e-classes.
% it searches for all $f$-application e-nodes $n_f$, adds $\alpha\mapsto
% n_f.\textit{child}_1$ to substitution $\sigma$, iterates through all
% $g$-application e-nodes $n_g$ in e-class $n_f.\textit{child}_2$, and only
% yield $\sigma$ if $n_g.\textit{child}_1=\sigma(\alpha)$. In general, this
% procedure runs in time that is quadratic of the \egraph~size. In a large
% e-graph, there may be thousands of pairs of $n_f$ and $n_g$ where $n_g$ is in
% e-class $n_f.\textit{child}_2$, but only a few satisfy the constraint
% $n_f.\textit{child}_1=n_g.\textit{child}_1$.

This inefficiency is due to the fact that na\"ive backtracking does not use the
equality constraints to prune the search space \textit{globally}. Specifically,
the above \ematching~pattern corresponds to three constraints for a potential
matching term $t$:
\begin{enumerate}
\item $t$ should have function symbol $f$.
\item $t$'s second child should have function symbol $g$.
\item $t$'s first child should be equivalent
      to the child of $t$'s second child.
\end{enumerate}

We can categorize these constraints into two kinds:
\begin{itemize}
  \item \textit{Structural constraints} are derived from the structure of the pattern.
    The structure of pattern
    $f(\alpha, g(\alpha))$ constrains the root symbol and the second symbol
    to be $f$ and $g$ respectively (i.e., constraints 1 and 2).
  \item \textit{Equality constraints}
    are implied by multiple occurrences of the same variable.
    Here, the occurrences of $\alpha$ implies that the terms at these positions
    should be equivalent with each other for all matches (i.e., constraint 3),
     which we call {\it equality constraints}.
    Following \citet{moskal}, we define patterns without equality constraints to be
    \textit{linear patterns}.
\end{itemize}

Backtracking search exploits the structural constraints first and defers
 checking the equality constraints to the end.
In our example pattern $f(\alpha, g(\alpha))$,
 backtracking search enumerates all $f(i, g(j))$,
 regardless of whether $i$ and $j$ are equivalent,
 only to discard nonequivalent matches later.
Complex query patterns may involve many variables
 that occur at several places, which will makes na\"ive backtracking search
 enumerate a very large number of candidates,
 even though the result size is small.
% In contrast, our approach exploits all equality constraints during query planning,
%  and can yield asymptotically better performance for non-linear patterns
%  and guarantees worst-case optimal run time in general.

\def\match{{\sf match}} \def\class{{\sf class}} \def\find{{\sf find}}
\def\dom{{\sf dom}}
\begin{figure}
  \begin{align*}
    \match(x,c,S) = & \{ \sigma \cup \{ x \mapsto c\} \mid \sigma \in S, x \not \in \dom(\beta)\}\ \cup\\
                    & \{ \sigma \mid \sigma \in S, \sigma(x) = c \}\\
                    % \match(c,t,S) = & c \in \class(t) \text{ ? } S \text{ : } \emptyset \\
    \match(f(p_{1}, \dots, p_{k}), c, S) = & \bigcup_{f(c_{1},\dots,c_{k})\in c}
                                             \match(p_{k}, c_{k}, \dots, \match(p_{1}, c_{1}, S))
  \end{align*}
  \caption{A declarative backtracking-based \ematching{} algorithm (reproduced
    from~\citet{efficient-ematching}). The set of substitutions for pattern $p$
    on \egraph~$G$ with \eclasses $C$ can be obtained by computing $\bigcup_{c\in C}\match(p, c, \emptyset)$.
    % \textbf{\color{red} maybe lets use an algo presentation to match up with GJ}
  }\label{alg:backtrack}
\end{figure}

\subsection{Conjunctive Queries}

Conjunctive queries are a subset of queries in relational algebra
 that use only select, project, and join operators
(as opposed to union, difference, or aggregation).
Conjunctive queries have many desirable theoretical properties
 (like computable equivalence checking),
 and they enjoy efficient execution
 thanks to decades of research from the databases community.

\subsubsection*{Relational Databases.}
A relational schema $S_D$ over domain $D$ is a set of relation symbols with
 associated arities.
A relation $R$ under a schema $S_D$ is a set of tuples;
 for each tuple $(t_1,\ldots,t_k) \in R$,
 $k$ is the arity of $R$ in $S_D$ and $t_i$ is an element in $D$.
A database instance (or simply database) $I$ of $S_D$ is a set of relations under $S_D$.

We use the notation $R(x, y).x$ to denote projection,
 i.e., $R(x, y).x = \{x \mid (x, y) \in R \}$.

\subsubsection*{Conjunctive Queries.}
A conjunctive query $Q$ over the schema $S_D$ is a formula of the form:
\[
  Q(x_1,\ldots x_k) \gets R_1(x_{1,1},\ldots,x_{1,k_1}),\ldots,
  R_n(x_{n,1},\ldots,x_{n,k_n}),
\]
where $R_1\ldots R_n$ are relation symbols in $S_D$ with arities
$k_1,\ldots k_n$ and the $x$ are variables.\footnote{
  Some definitions of conjunctive queries allow both variables and constants.
  We only allow variables without loss of generality:
  any constant $c$ can be specified with a distinguished relation $R_{c} = \{c\}$.}
We call the $Q(\ldots)$ part the \textit{head} of the query,
 the remainder is the \textit{body}.
Each $R_i(\ldots)$ is called an {\em atom}.
Variables that appear in the head are called \textit{free variables},
 and they must appear in the body.
Variables that appear in the body but not the head are called \textit{bound variables},
 since they are implicitly existentially quantified.

% $v_{i,j}$ are variables in $V$.
% and $x_1,\ldots x_k$ are variables
% occurring in $v_{i,j}$.

\subsubsection*{Semantics of Conjunctive Queries.}
Similar to \ematching,
 evaluating a conjunctive query $Q$ yields substitutions.
Specifically,
 evaluation yields substitutions
 that map free variables in $Q$ to
 elements in the domain such that
 there exists a mapping of the bound variables
 that causes every substituted atom to be present in the database.
Bound variables are {\em projected out} and not present in
 resulting the substitutions.

% \textbf{\color{blue} This definition may cause problems if variables in
%   the body do not occur in the head.
% Instead, we may say the result of CQ is a set of tuples, instantiated from a
% substitution (of all variables in the body) such
% that ... Example:
% [...] Evaluating $Q$ over $I$ yield the set of all possible
%  $(\sigma(x_1), \ldots, \sigma(x_n))$,
%  where $\sigma$ is a substitution satisfying
%  \[
%    \forall j\in [n].\
%    R_j\left( \sigma(x_{j,1}),\ldots, \sigma(x_{j,k_j})\right) \in I.
%  \]
% }

More formally,
 let $I$ be a database of schema $S_D$ and
 let $Q$ be a conjunctive query over the same schema
 with $k$ variables in its head.
Let the $n$ atoms in the body of
 $Q$ be $R_1, \ldots, R_n$
 where $R_j$ has arity $k_j$.
Evaluating $Q$ over $I$ yields
 a substitution
 $\sigma = \{ x_1 \mapsto t_1, \ldots, x_k \mapsto t_k \}$
 iff
 there exists a $\sigma' \supset \sigma$ mapping
 all variables in $Q$ such that:
\[
  \bigwedge_{j \in [n]}(\sigma'(x_{j,1}),\ldots,\sigma'(x_{j,k_j})) \in R_{j}
\]
% Evaluating $Q$ over $I$ yields substitutions
%  % $\sigma = \{ x_1 \mapsto t_1, \ldots, x_k \mapsto t_k \}$
%  $\sigma = \{ x_1 \mapsto \sigma'(x_1), \ldots, x_k \mapsto \sigma'(x_k) \}$,
%  where $\sigma'$ is a substitution subsuming $\sigma$
%  that map variables from the bodies of $Q$ to elements from $D$.
% Specifically, $\sigma'$ satisfies the following:
% \[
%   \forall j \in [n].\
%   \left( \sigma'(x_{j,1}), \ldots, \sigma'(x_{j,k_j}) \right) \in R_j
% \]
% In other words, $Q$ returns the set (where $x_{l}, \ldots, x_{m}$ are body variables not in the head):
%
%
%  $\sigma = \{ x_1 \maps t_1, \ldots, x_k \maps t_k \}$
%  that map all $k$ variables
%  in the head of $Q$ to elements of $D$ such that

% $R_i(\sigma(t_1),\ldots \sigma(t_k))$ are atoms in
%  $(\sigma(x_1),\ldots,\sigma(x_k))$ where $\sigma$ is a
%  substitution satisfying that $R_i(\sigma(t_1),\ldots \sigma(t_k))$ are atoms in
% $I$ for $i=1,\ldots,n$.
% We denote the result set as $Q(I)$.

% Similar to \ematching~substitutions, we define a conjunctive query substitution
% as a function that maps every variable occurring in it to elements in
% $D$. The semantics of conjunctive queries is defined as follows: let $Q$ be a
% conjunctive query, $I$ be an database of schema of $S_D$.
% Conjunctive query $Q$ yields all tuples
% set of all $\textit{ans}(\sigma(x_1),\ldots,\sigma(x_k))$ where $\sigma$ is a
% substitution satisfying that $R_i(\sigma(t_1),\ldots \sigma(t_k))$ are atoms in
% $I$ for $i=1,\ldots,n$. We denote the result set as $q(I)$.

In practice, conjunctive queries are often evaluated according to a \textit{query plan}
 which dictates each step of execution.
For example,
 many industrial database systems will construct tree-like query plans,
 where each node describes an operation like scanning a relation or joining two
 intermediate relations.
Industrial database systems typically construct query plans
 based on binary join algorithms such as hash joins and merge-sort join,
 which process two relations at a time.
The quality of a query plan critically
 determines the performance of evaluating a conjunctive query.

We observe that conjunctive query and \ematching
 are structurally similar: both are defined as finding
 substitutions whose instantiations are present in a database.
% In fact, the \ematching problem can be viewed as ``nested''
%  conjunctive queries on a database where the relations are nested.
Therefore, it is tempting to reduce \ematching to
 a conjunctive query over the relational database,
 thereby benefiting from well-studied techniques from the databases community,
 including join algorithms and query optimization.
We achieve exactly this in \autoref{sec:algorithm}.

\subsection{Worst-Case Optimal Join Algorithms}

% \begin{figure}
%   \begin{subfigure}{0.45\textwidth}
%     \centering \tikzfig{worstcase}
%     \caption{ Input relations. }\label{fig:worst}
%   \end{subfigure}
%   \begin{subfigure}{0.45\textwidth}
%     \centering \tikzfig{triangle}
%     \caption{ $Q \leftarrow R(x,y), S(y,z), T(z, x)$. }\label{fig:triangle}
%   \end{subfigure}
%   \caption{ Example input and query hypergraph of $Q$. }\label{fig:wcoj}
% \end{figure}

% \begin{figure}
% \begin{align*}
% \textsf{GJ}(\bigwedge_i R_i(\overline{x_i}), [x_1] \concat \overline{x}, \sigma) & =
% \{ \textsf{GJ}(x_1 = a \wedge \bigwedge_i R_i(\overline{x_i}), \overline{x}, \sigma \cup \{ x_{1} \mapsto a \} ) \mid a \in \bigcap_{x_1 \in R_j} R_j \} \\
% \textsf{GJ}(\bigwedge_i R_i(\overline{x_i}), \emptyset, \sigma) & = \sigma
% \end{align*}
% \caption{Generic Join algorithm. \textsf{GJ} takes 3 arguments:
%  a conjunction of atoms $\bigwedge_i R_i(\overline{x_i})$,
%  a list of variables indicating the variable ordering,
%  and a (partial) substitution $\sigma$.}
% \end{figure}

The run time of any algorithm for answering conjunctive queries
 is lower-bounded by the output size,
 assuming the output must be materialized.
How large can the output of a conjunctive query be on a particular database?
A na{\"i}ve bound is simply the product of the size of each relation,
 which is the size of their cartesian product.
Such a na{\"i}ve bound fails to consider the query's structure.
The AGM bound~\cite{agm}
 gives us a bound for the worst-case output size.
In fact, the AGM bound is tight;
 there always exists a database
 where the query output size is the AGM bound.

The AGM bound and worst-case optimal joins
 are recent developments in databases research.
We do not attempt to provide a comprehensive
 background on these topics here;
 familiarity with the existence of the AGM bound
 and the generic join algorithm is sufficient for this paper.

Consider $Q(x,y,z) \leftarrow R(x,y), S(y,z), T(x,z) $,
 also known as the ``triangle query'',
 since output tuples are triangles between the edge relations $R,S,T$.
We calculate a trivial bound $|Q| \leq |R| \times |S| \times |T|$.
If $|R|=|S|=|T|=N$, then $|Q| \leq N^3$.
We can derive a tighter bound from $|Q| \leq |R| \times |S| = N^2$.
That is because $Q$
contains fewer tuples than the query $ Q’(x,y,z) \leftarrow R(x,y), S(y,z) $
as $Q$ further requires $(x,z) \in T$.
The AGM bound for $Q$ is even smaller: $N^{3/2}$.
It is computed from the \emph{fractional edge
  cover} of the \emph{query hypergraph}.

\subsubsection*{Query Hypergraph.}
\label{sec:hypergraph}

\begin{figure}
  \centering \tikzfig{triangle}
  \caption{ Query hypergraph of $Q(x, y, z) \gets R(x,y), S(y,z), T(z, x)$. }\label{fig:wcoj}
\end{figure}

The hypergraph of a query is simply the hypergraph
 with
 a vertex for each variable
 and a (hyper)edge for each atom.
The edge for an atom $R(x_{i}, \ldots, x_{k})$ connects the vertices
 corresponding to the variables $x_{i}, \ldots, x_{k}$.
\autoref{fig:wcoj} illustrates $Q$'s hypergraph.
% Certain hypergraphs are called \emph{cyclic}.
% The definition of a hypergraph cycle is related to but not exactly
%  that of cycles in conventional graphs.

\subsubsection*{Cyclic and Acyclic Queries.}
Certain queries can be represented by a tree,
 called the \textit{join tree},
 where each node corresponds to an atom in the query.
Furthermore, for each variable $x$ the nodes
 corresponding to the atom containing $x$
 must form a connected subtree.
Queries that admit such a join tree
 are said to be \textit{acyclic};
 otherwise, the query is \textit{cyclic}.\footnote{
  A cycle in the hypergraph does not necessarily entail a cyclic query,
   since the hypergraph may still admit a join tree.
}
The triangle query is cyclic because it cannot
 be represented by a join tree.
Acyclic queries can be answered more efficiently
 than cyclic ones.

\subsubsection*{Fractional Edge Cover.}
A set of edges \emph{cover} a graph if they touch all vertices. For $Q$'s
hypergraph, any two edges form a cover.
A \emph{fractional edge cover} assigns a
 weight in the interval $[0, 1]$ to each edge
 such that,
 for each vertex $v$,
 the weights of the edges containing $v$ sum to at least 1.
Every edge cover is a fractional cover,
 where every edge is assigned a weight of 1 if it is in the cover,
 and 0 otherwise.
For $Q$'s hypergraph,
 $\{R \mapsto 1/2, S \mapsto 1/2, T \mapsto 1/2\}$
 is the fractional edge cover with lowest total weight.

\subsubsection*{The AGM Bound.}
The AGM bound~\cite{agm} for a query with body atoms
 $R_{i}(\ldots)$ for $i \in [k]$
 is defined as $\min_{w_1,\ldots,w_k} \Pi_{i \in [k]} |R_{i}|^{w_i}$,
 where $\{R_{i} \mapsto w_{i} \mid i \in [k] \}$ forms a fractional edge cover.
For example, the AGM bound for $Q$ is
 $|R|^{1/2}|S|^{1/2}|T|^{1/2} = N^{3/2}$ when $|R| = |S| = |T| = N$.
This is the upper bound of $Q$’s output size;
 i.e.\ in the worst case $Q$ outputs this many tuples.

\subsubsection*{Generic Join.}

A desirable algorithm for answering conjunctive queries
 should run in time linear to the worst case output size.
Recent developments in the databases community have led to
 such an algorithm~\cite{wcoj}, one of which is \textit{generic join}~\cite{gj}.
Generic join has one parameter:
 an ordering of the variables in the query.
Any ordering guarantees a run time linear to the worst-case output size,
 but different orderings can lead to dramatically different run time
 in practice~\cite{emptyheaded}.

\begin{algorithm}
  \LinesNumbered{} \SetAlgoLined{}
  \KwResult{$\text{GJ}(Q, \emptyset)$ computes the output of query $Q$}
  \KwIn{query $Q$, partial substitution $\sigma$}
  \tcc{$k$ indicates how many variables remain in query $Q$}
  \eIf(\tcc*[f]{there are no more variables, so $\sigma$ is complete}){$k=0$}
    {output $\sigma$}
    {
      choose a variable $x$\;
      \tcc{
        Compute $D_x$, which all possible values of $x$, by intersecting the attributes of the relations where $x$ occurs.
        Intersection must be computed in $O(\min(|R_j.x|))$ time.}
      $J = \{j \mid x \in \overline{X}_j\}$\;
      $D_x = \cap_{j \in J} R_j.x$\label{ln:gj-isect}\;
      \For()
      {$v \in D_x$}
      {
        \tcc{compute residual query by replacing variable $x$ with constant $v$}
        $Q' = Q[v/x]$\;
        GJ$(Q', \sigma \cup \{ x \mapsto v \})$
      }
    }
  \caption{
    Generic join for the general query
    $Q(x_1,\ldots,x_k) \gets R_1(\overline{X}_1),\ldots,R_n(\overline{X}_n)$
  }\label{alg:gj-general}
\end{algorithm}

\autoref{alg:gj-general} shows the generic join algorithm.
Generic join is recursive,
 proceeding in two steps.
First,
 it chooses a variable from the query and
 collects all possible values for that variable in the query.
Then,
 for each of those values,
 it builds a \textit{residual query}
 by replacing occurrences of the variable in the query with a possible
 value for that variable.
These residual queries are solved recursively,
 and when there are no more variables in the residual query,
 the algorithm yields the substitution it has accumulated so far.

Generic join
 requires two important performance
 bounds to be met in order for its own run time to meet the AGM bound.
First, the intersection on line \ref{ln:gj-isect}
 must run in $O(\min(|R_j.x|))$ time.
Second,
 the residual relations
 should be computed in constant time,
 i.e., computing from the relation $R(x, y)$
 the relation $R(v_x, y)$ for some $v_x \in R(x, y).x$
 must take constant time.
Both of these
 can be solved by using tries
 (sometimes called prefix or suffix trees)
 as an indexing data structure.
\autoref{fig:simple-trie} shows an example trie.
Tries allow fast
 intersection because
 each node is a map which
 can be intersected in time linear to the size of the smaller map.
Tries also allow constant-time access
 to residual relations according to a compatible variable ordering.

 \begin{figure}
   \begin{subfigure}[b]{0.45\linewidth}
     \centering
     \begin{tabular}{|lll|}
       \hline
       \bf x & \bf y & \bf z \\
       \hline
       1 & 2 & 4 \\
       1 & 2 & 6 \\
       1 & 3 & 7 \\
       8 & 2 & 4 \\
       \hline
     \end{tabular}
     \caption{Table for relation $R(x, y, z)$.}
   \end{subfigure}
   \hfill
   \begin{subfigure}[b]{0.5\linewidth}
     \centering
     \includegraphics[height=22mm]{simple-trie.pdf}
     \caption{Trie for relation $R(x, y, z)$ using ordering $[x, y, z]$.}
   \end{subfigure}
   \caption{
     A trie is a tree where
       every node is a map (typically a hashmap or sorted map)
       from a value to a trie.
     Every path from the root of a trie to a leaf represents a tuple in the relation.
     Tries allow efficient computation of residual relations.
     For example, $R(1, y, z)$ can be computed quickly by following the
       $x=1$ edge from the root.
   }
   \label{fig:simple-trie}
 \end{figure}


A useful way to understand generic join is to observe the
 loops and intersections it performs on a specific query.
\autoref{alg:gj} shows generic join computing the triangle query.
Given a variable ordering ($[x,y,z]$ in this case),
 generic join assumes the input relations are stored in tries according to the ordering,
 so $R(x,y)$ is
 stored in a trie with $x$s on the first level, and $y$s on the second.
This makes accessing the residual relations ($R(v_x, y)$) fast since
 the replacement of variables with values
 is done according to the given variable ordering.
Note how the algorithm is essentially just nested \texttt{for} loops.
There is no explicit filtering step;
 the intersection of residual queries guarantees
 that once a complete tuple of values is selected,
 it can be immediately output without additional checking.

\begin{algorithm}
  \LinesNumbered{} \SetAlgoLined{}
  \KwResult{compute $Q(x,y,z) \gets R(x,y),S(y,z),T(z,x)$}
  $X = R(x,y).x \cap T(z,x).x$\;
  \For(\tcc*[f]{compute $Q(v_x,y,z) = R(v_x,y),S(y,z),T(z,v_x)$})
  {$v_x \in X$}
  {
    $Y = R(v_x,y).y \cap S(y,z).y$\;
    \For(\tcc*[f]{compute $Q(v_x,v_y,z) = R(v_x,v_y),S(v_y,z),T(z,v_x)$})
    {$v_y \in Y$}
    {
      $Z = S(v_y,z).z \cap T(z,v_x).z$\;
      \For(\tcc*[f]{yield join results $Q(v_x, v_y, v_z)$})
      {$v_z \in Z$}
      {output$(v_x,v_y,v_z)$ } } }
  \caption{
    Generic join for the triangle query, with ordering $[x, y, z]$.
  }\label{alg:gj}
\end{algorithm}

% \subsubsection{When to Use Generic Join.}

% Use Generic Join when the query is cyclic - if the query is acyclic, the
% theoretical optimal algorithm is Yannakakis (equivalent to message passing in
% probabilistic graphical models) which runs in linear time in the input size;
% in practice people use binary joins for acyclic queries, because Yannakakis
% always incurs a constant factor of 3 (it’s a 3-pass algorithm).

% In practice, Generic Join may have good performance compared to binary joins,
% since it is sometimes equivalent to a linear join tree. The “average case”
% complexity of Generic Join is open.


%%% Local Variables:
%%% mode: latex
%%% TeX-master: "main"
%%% End:
